Сначала оценивающие последовательности собираются в \hyperref[note:C038105E-BD42-44A8-9542-2DBE9FFAAA3F]{теорему об ускоренном методе Нестерова}: множитель $1 - \sqrt{\mu/L}$ вместо $1 - \mu/L$ у градиентного спуска. Близкую скорость локально даёт \hyperref[note:C6A43F85-13E0-46B3-96A8-BCB9F34ECD57]{метод тяжёлого шарика}. Затем глава возвращается к Ньютону: \hyperref[note:6BA18878-DE34-46A5-ABE2-EDA09CAE56A7]{квазиньютоновские методы} заменяют гессиан матрицей $B_k$, которая пересчитывается дёшево и подчиняется \hyperref[note:E976DE49-0091-45E1-B621-A236F8B87859]{условию секущей}; примеры --- BFGS и SR1.
